\documentclass[11pt]{article}
\usepackage{amsmath,amssymb}

\title{Minimum Makespan Scheduling}
\author{}
\date{}

\begin{document}
\maketitle

\section*{Problem}
Given the times taken to complete $n$ jobs, $t_1,t_2,\dots,t_n$, and $m$ identical machines, assign each job to a machine so that the total time taken to finish all jobs is minimized. For example, if $n=5$, $m=3$, and the job times are $5,4,4,6,7$ hours, the best schedule assigns jobs taking $4$ and $5$ hours to machine 1, jobs taking $4$ and $6$ hours to machine 2, and the $7$-hour job to machine 3. The total time is then $10$ hours.

Consider the following greedy algorithm: order the jobs arbitrarily and, in this order, assign each job to the machine that has been given the least work so far. Let $T_{\mathrm{OPT}}$ be the total time taken by an optimal schedule, and let $T_A$ be the time taken by the greedy algorithm. Show that
\[
  \frac{T_A}{T_{\mathrm{OPT}}} \leq 2.
\]
In other words, this algorithm always finds a schedule taking at most twice as long as an optimal schedule.

\section*{Solution}
Without loss of generality, let the arbitrary ordering be $t_1,t_2,\dots,t_n$. We have
\[
  T_{\mathrm{OPT}}
  \geq \max\left\{\frac{\sum_{i=1}^n t_i}{m}, \max\{t_1,t_2,\dots,t_n\}\right\}.
\]
Therefore,
\[
  2T_{\mathrm{OPT}}
  \geq \frac{\sum_{i=1}^n t_i}{m} + \max\{t_1,t_2,\dots,t_n\}.
\]

Now, without loss of generality, let machine 1 be a machine with maximum load, so
\[
  T_A=\sum_{i\in \mathrm{machine}_1} t_i.
\]
Let $l$ be the last job assigned to machine 1. Immediately before assigning job $l$, machine 1 had the minimum load among all $m$ machines. Its load at that time is therefore at most the average load,
\[
  \frac{\sum_{i=1}^n t_i}{m}.
\]
Consequently,
\[
  T_A \leq \frac{\sum_{i=1}^n t_i}{m}+t_l
  \leq \frac{\sum_{i=1}^n t_i}{m}+\max\{t_1,t_2,\dots,t_n\}
  \leq 2T_{\mathrm{OPT}}.
\]
Thus $\frac{T_A}{T_{\mathrm{OPT}}}\leq 2$. \qed

\end{document}
